\documentclass[UTF8,a4paper,11pt,twoside]{ctexrep}
\usepackage[margin=2.3cm]{geometry}
\usepackage{amsmath,amssymb}
\usepackage{booktabs,longtable,array,tabularx}
\usepackage{listings,xcolor}
\usepackage{tikz}
\usetikzlibrary{arrows.meta,positioning,calc,fit,shapes.geometric,
                decorations.pathreplacing,decorations.markings,patterns}
\usepackage{hyperref}
\hypersetup{colorlinks=true,linkcolor=blue!60!black,urlcolor=blue!60!black,
            citecolor=blue!60!black}
\usepackage{enumitem}
\setlist{nosep,leftmargin=1.6em}

% --- 代码样式 ---
\lstset{
  basicstyle=\ttfamily\footnotesize,
  breaklines=true,columns=fullflexible,
  frame=single,framerule=0.3pt,rulecolor=\color{gray!50},
  backgroundcolor=\color{gray!8},
  keywordstyle=\color{blue!60!black}\bfseries,
  commentstyle=\color{green!50!black},
  stringstyle=\color{red!60!black},
  language=C,
  morekeywords={module,endmodule,analog,electrical,input,output,inout,
                parameter,from,integer,real,initial_step,cross,transition,
                begin,end,case,endcase,if,else,for,posedge,negedge},
}

% --- TikZ 快捷样式 ---
\tikzset{
  blk/.style={draw,rounded corners=2pt,minimum height=7mm,minimum width=14mm,
              font=\scriptsize,align=center,fill=white},
  sig/.style={-{Stealth[length=2mm]},thick},
  dashsig/.style={-{Stealth[length=2mm]},thick,dashed},
}

% --- 标题信息 ---
\title{\bfseries 12 位冗余 SAR ADC 前台校准\\[2pt]
       \large 行为级实现、算法验证与性能分析}
\author{SAR\_ADC\_IDEAL\_CAL\_REPLACEMENT\_V6\_CLEAN}
\date{2026 年 7 月}

\begin{document}
\maketitle

%===================================================================
\begin{abstract}
\noindent
本报告描述一套 12 位非二进制冗余 SAR ADC 的 Verilog-A 行为级替代包，
用于验证 Huang~2024 前台位权自校准算法。包含 5 个核心模块与 2 套蒙特卡洛
CDAC 模型，实现了 offset-binary SAR 搜索、D+/D$-$ 差分失调消除、32 对 Q 格式
平均、确定性 sub-LSB dither 注入等完整算法链。理论校准分辨率达
$1.1\times10^{-6}$，实测 ENOB 11.5--11.7 位。报告定量分析了 ENOB 损耗的四大
根因（递归误差传播、冗余裕量过紧、端点裁剪、dither-SAR 耦合），给出四级改进
路径，并与 Huang~2024 论文的 16-bit 晶体管级实现做逐维度对比。结论：当前
ENOB 瓶颈在架构而非校准精度；对中小失配，DEM 的 $\sqrt{N}$ 效应在理论上优于
前台校准，但需 CDAC 硬件重设计。
\end{abstract}

\tableofcontents

%===================================================================
\chapter{引言}
%===================================================================

\section{问题定义}

高分辨率 SAR ADC 的线性度受限于 CDAC 电容失配。对于 12 位 ADC，要求电容匹配
精度优于 $0.024\%$，远超标准工艺的 $3\sigma$ 失配水平（$\sim 0.1\%$）。
解决途径有二：\textbf{校准}（测量位权并修正解码）与\textbf{动态元件匹配}（轮换
元件将失配转为噪声）。Huang~2024 论文\cite{huang2024}提出一套基于 offset-binary
SAR 搜索与 D+/D$-$ 差分的前台位权自校准方案，在 16-bit 5-MS/s 原型上实现
SNDR 93.7\,dB。

本项目将 Huang~2024 算法移植到 12 位非二进制冗余 SAR ADC 的 Verilog-A 行为级
平台，目标是在晶体管级数字链路尚未就绪时验证校准算法的正确性与边界。

\section{设计约束}

\begin{itemize}
\item CDAC 为非二进制权重：$2048,1024,512,256,256R,128,48,32,20,12,8,4,2$，
      总权重 4351，冗余位 $256R$ 提供纠错重叠
\item calDAC 复用 CDAC 低位段（权重 48--2，7 位），offset-binary 编码
\item 仅校准 3 位 MSB（$w_2/w_1/w_0$，权重 512/1024/2048），中位裸奔
\item 理想比较器无噪声，需确定性 dither 制造亚 LSB 可观测性
\end{itemize}

%===================================================================
\chapter{系统架构}
%===================================================================

\section{顶层信号流}

图~\ref{fig:topflow} 展示五个 Verilog-A 模块的信号流。外部 \texttt{CLK00}
启动 \texttt{SYNC\_asnyc} 产生 14 个比较器周期；\texttt{DEC\_CAL\_PHY} 作为
校准引擎与解码器的合一模块，\textbf{是最终比较器时钟 \texttt{CLK} 的唯一驱动者}
——这一约束消除了 SYNC 与 DEC 对 CLK 网络的双驱动冲突。

\begin{figure}[ht]\centering
\begin{tikzpicture}[node distance=0.6cm and 1.2cm]
\node[blk,fill=blue!10] (sync) {SYNC\_asnyc\\定时发生};
\node[blk,right=2.2cm of sync,fill=green!10] (sar) {SAR\_LOGIC\\14拍逻辑};
\node[blk,below=1.0cm of sar,fill=orange!10] (sw) {SWITCH\_CAL\\CDAC开关};
\node[blk,below=1.0cm of sync,fill=red!10] (com) {COM\_ideal\\比较器};
\node[blk,left=1.2cm of com,fill=purple!10] (dec) {DEC\_CAL\_PHY\\校准+解码};

\draw[sig] (sync) -- node[above,font=\scriptsize]{CLK\_SAR} (sar);
\draw[sig] (dec) -- node[above,font=\scriptsize]{CLK} (com);
\draw[dashsig] (dec.south) to[bend right=20] node[below,font=\scriptsize]{CAL} (com.south);
\draw[sig] (sar) -- node[right,font=\scriptsize]{BP$[13{:}0]$} (sw);
\draw[sig] (com) -- node[right,font=\scriptsize]{COMP/COMN} (sar.west);
\draw[dashsig] (dec.east) to[bend left=15] node[above,font=\scriptsize]{BITD\_CAL} (sw.west);
\draw[sig] (sw.east) to[bend right=10] ++(1.5,0) node[right,font=\scriptsize]{顶板电压};
\node[font=\scriptsize,above=0.1cm of com] {$V_P,V_N$};
\draw[sig] (dec) -- ++(0,-1.3) node[below,font=\scriptsize]{Bit$[11{:}0]$};
\end{tikzpicture}
\caption{顶层信号流。实线为时钟/数据，虚线为校准控制。}
\label{fig:topflow}
\end{figure}

\section{模块职责与关键参数}

\footnotesize
\begin{longtable}{p{3.2cm}p{2.0cm}p{8.5cm}}
\toprule
\textbf{模块} & \textbf{角色} & \textbf{关键参数与职责} \\
\hline
\endhead
\texttt{COM\_ideal} & 比较器 & \texttt{CAL\_LSB\_V}=0.8274\,mV, \texttt{DITHER\_LEVELS}=32；CAL=1 注入 dither，CAL=0 严格为零 \\
\texttt{SWITCH\_CAL} & CDAC 开关 & \texttt{REDUNDANT\_SAMPLES\_VCM}=1（满量程修复）；三态编码 10/01/00 \\
\texttt{SAR\_LOGIC} & SAR 逻辑 & 14 拍 compare-then-commit；\texttt{NORMAL\_REF\_SWAP}=1 \\
\texttt{SYNC\_asnyc} & 定时 & 4 相时钟，4\,ns/决策，14 决策共 56\,ns \\
\texttt{DEC\_CAL\_PHY} & 校准+解码 & \texttt{AVG\_PAIRS\_LOG2}=5, \texttt{Q\_SCALE}=64, \texttt{WEIGHT\_TOL}=0.35；\texttt{ENFORCE\_MIN\_AVG}=1（合并版） \\
\bottomrule
\end{longtable}

%===================================================================
\chapter{CDAC 架构与满量程修复}
%===================================================================

\section{非二进制权重设计}

CDAC 物理权重序列在 \texttt{SWITCH\_CAL.va} 中定义：
\[
\underbrace{2048,1024,512}_{\text{MSB 段}},\ 
\underbrace{256,256R}_{\text{冗余段}},\ 
\underbrace{128,48,32,20,12,8,4,2}_{\text{calDAC 段}},\ 
\underbrace{1}_{\text{终端}}
\]
总物理权重 $W_{total}=4351$。冗余位 $256R$ 在 stage-3 与 stage-4 间提供
$256$ 的纠错重叠范围。calDAC 段（权重 $2$--$48$，7 位）在正常转换时作为
普通低位电容，在校准模式复用为测量 DAC。

\section{满量程修复}

\textbf{问题}：若冗余 $256R$ 支路与所有信号电容一起采样 $V_{IP}/V_{IN}$，则
采样信号跨度 $=4351$ 间隔。将 4351 个间隔压缩到 4095 个输出码会周期性合并
相邻码字，在 FFT 中表现为谐波，将理想 SNDR 上限限制在 $\approx 71$\,dB。

\textbf{修复方案}（\texttt{SWITCH\_CAL.va} L10--L22）：令冗余支路在采样相
接 $V_{CM}$ 而非 $V_{IP}/V_{IN}$，使其不贡献信号电荷。转换与校准相仍切换到
$V_{REFP}/V_{REFN}$，保留纠错能力。修复后信号权重和 $=4094$，加终端比较器
决策的 $1$，恰好 $4095$ 间隔，恢复 12 位满量程。

\begin{figure}[ht]\centering
\begin{tikzpicture}[font=\scriptsize]
% 信号电容
\foreach \i/\w/\lab in {0/2.0/2048,1/1.4/1024,2/1.0/512,3/0.7/256,4/0.7/256R,5/0.5/128,6/0.4/48,7/0.3/32,8/0.25/20,9/0.2/12,10/0.15/8,11/0.12/4,12/0.1/2}{
  \draw[fill=blue!20] (\i*0.9,0) rectangle ({\i*0.9+0.7},{\w*0.3});
  \node[font=\scriptsize] at ({\i*0.9+0.35},{\w*0.3+0.15}) {\lab};
}
\draw[fill=red!30] (4*0.9,0) rectangle ({4*0.9+0.7},{0.7*0.3});
\node[font=\scriptsize,red] at ({4*0.9+0.35},{0.7*0.3+0.15}) {256R};
\draw[<->] (0,-0.2) -- ({12*0.9+0.7},-0.2) node[midway,below]{信号电容};
\draw[<->] ({4*0.9-0.1},-0.2) -- ({4*0.9+0.8},-0.2) node[midway,below,red,font=\scriptsize]{采 VCM};
\node[right] at ({12*0.9+0.9},0.3) {$\sum=4094+1=4095$};
\end{tikzpicture}
\caption{CDAC 权重分布与冗余支路采样修复。红色 $256R$ 采样 $V_{CM}$，不贡献信号。}
\label{fig:cdac}
\end{figure}

\section{三态开关编码}

每个电容 $k$（转换级索引 $=13-k$）在三个相位下的控制：

\begin{table}[ht]\centering\footnotesize
\begin{tabular}{lccl}
\toprule
\textbf{相位} & \textbf{信号电容} & \textbf{冗余支路} & \textbf{控制信号} \\
\hline
采样 & $V_{IP}/V_{IN}$ & $V_{CM}$ & BITD[k], BITU[k] \\
转换 & $10\to V_{REFP}$, $01\to V_{REFN}$ & 同左 & SET + BITP/BITN \\
校准 & $10=+$, $01=-$, $00=$VCM & 同左 & BITD\_CAL, BITU\_CAL \\
\bottomrule
\end{tabular}
\end{table}

%===================================================================
\chapter{前台校准算法}
%===================================================================

\section{递归校准原理}

校准从低位 MSB 向高位递归推进。每次测量一个目标电容 $W_{target}$ 相对于
"墙"（已校准低位之和 $W_{wall}$）的权重差 $R$：

\begin{table}[ht]\centering\footnotesize
\begin{tabular}{cccccl}
\toprule
\textbf{序} & \textbf{目标} & \textbf{权重} & \textbf{墙} & \textbf{墙权重} & \textbf{测量结果} \\
\hline
1 & C10 & 512 & $w_3+w_4$ & 512 & $w_2$ \\
2 & C11 & 1024 & $w_2+w_3+w_4$ & 1024 & $w_1$ \\
3 & C12 & 2048 & $w_1+w_2+w_3+w_4$ & 2048 & $w_0$ \\
\bottomrule
\end{tabular}
\end{table}

递归性体现在：$w_2$ 的测量结果进入 $w_1$ 的墙，$w_1$ 的结果再进入 $w_0$ 的墙。
因此 $w_2$ 的测量误差会\textbf{传播}到 $w_1$ 和 $w_0$。

\section{offset-binary SAR 搜索}

calDAC 初始状态为全 10（输出 $+126$），随后逐位从 10 切到 01，使输出减少
$2\times$ weight。试探方向\textbf{反向于残差} $R$：若 $R>0$（目标偏大），
calDAC 输出向负方向搜索；若 $R<0$，向正方向搜索。7 步搜索后 calDAC 收敛，
得到 0--127 的码 $D$，加上终端比较器决策 $b_T$：
\[
D = \sum_{\text{switched}} w_i + b_T
\]

\section{D+/D$-$ 差分消除比较器失调}

对同一目标做正反两次测量，残差符号翻转，比较器失调 $V_{os}$ 同号：
\begin{align}
D_+ &\approx \frac{W_{target}-W_{wall}+V_{os}}{\Delta} \\
D_- &\approx \frac{-W_{target}+W_{wall}+V_{os}}{\Delta}
\end{align}
差分后失调一阶消除：
\[
\boxed{R = D_+ - D_- \approx \frac{2(W_{target}-W_{wall})}{\Delta}}
\]
本实现不除以 2，而是通过 Q 格式系数吸收因子 2。

\section{Q 格式平均}

32 对 $(D_+,D_-)$ 的差分结果做 Q 格式平均，保留亚 LSB 小数：
\[
R_q = \mathrm{round}\!\left(q\_scale \cdot \frac{1}{N}\sum_{i=1}^{N}(D_{+,i}-D_{-,i})\right),\quad
q\_scale=64,\ N=32
\]
测得权重 $W_{meas} = W_{wall,q} + R_q$，精度 6 位小数。容差门限 $\pm 35\%$
（\texttt{WEIGHT\_TOL}=0.35），超限则 REJECT。

\section{校准状态机}

\texttt{DEC\_CAL\_PHY.va} 定义 6 态状态机（图~\ref{fig:fsm}）：
CAL\_CLK 上升沿设置试探（L330--L375），下降沿锁存比较器并决策 keep/undo
（L378--L486），RST1 上升沿结束一个方向（L489--L586）。

\begin{figure}[ht]\centering
\begin{tikzpicture}[node distance=1.4cm and 1.0cm,font=\scriptsize]
\node[blk] (idle) {IDLE};
\node[blk,right=of idle] (pre) {PRECHARGE};
\node[blk,right=of pre] (trial) {SAR\_TRIAL};
\node[blk,below=1.0cm of trial] (term) {TERMINAL};
\node[blk,left=of term] (fin) {FINISH\_DIR};
\node[blk,left=of fin,fill=green!10] (done) {DONE};
\draw[sig] (idle)--(pre);
\draw[sig] (pre)--(trial);
\draw[sig] (trial)--(term);
\draw[sig] (term)--(fin);
\draw[sig] (fin)--(done);
\draw[dashsig] (fin.north) to[bend left=35] node[above,font=\scriptsize]{下一目标/方向} (trial.north);
\node[below=0.1cm of trial,font=\scriptsize,gray] {7步};
\node[below=0.1cm of term,font=\scriptsize,gray] {1步};
\end{tikzpicture}
\caption{校准状态机。SAR\_TRIAL 重复 7 次（7 位 calDAC），TERMINAL 1 次。}
\label{fig:fsm}
\end{figure}

\section{时钟多路选择}

\texttt{DEC\_CAL\_PHY.va} L729--L733 控制比较器时钟来源：
\[
\text{CLK} = \begin{cases}
\text{CLK\_SAR} & \text{cal\_mode}=0\ (\text{正常转换})\\
\text{CAL\_CLK} & \text{cal\_mode}=1\ \text{且 frame\_active}
\end{cases}
\]
校准模式下每个方向恰好 8 个比较器沿（7 步 SAR + 1 终端），16 沿对齐一个
D+/D$-$ 对。\texttt{SYNC\_asnyc} 的符号端 \texttt{CLK} 必须接到网络
\texttt{CLK\_SAR}，而非最终 \texttt{CLK}。

%===================================================================
\chapter{确定性 Dither 机制}
%===================================================================

\section{问题：理想比较器的可观测性缺失}

理想比较器无噪声，当残差 $R$ 落在 calDAC 两个量化电平之间时，SAR 永远收敛到
同一个整数码，\textbf{小数部分无法观测}。Huang~2024 论文依赖真实电路噪声
$\sigma_{n,tot}$ 提供亚 LSB 可观测性；本行为级平台必须用确定性 dither 替代。

\section{dither 原理}

dither 是一个\textbf{已知的、确定性的}阈值偏移，在 32 个测量对中均匀扫过
一个 calDAC LSB。对于残差 $R$ 的小数部分 $\{R\}$，不同 dither 电平会使终端
bit 翻转：

\begin{figure}[ht]\centering
\begin{tikzpicture}[font=\scriptsize]
\draw[->] (0,0)--(8,0) node[right]{dither 电平 $d$};
\draw[->] (0,0)--(0,2.5) node[above]{终端 bit $b_T$};
\draw[thick] (0,0.3)--(3.5,0.3)--(3.5,2.0)--(7,2.0);
\draw[dashed] (3.5,0)--(3.5,2.0);
\node[below] at (3.5,0) {$d=-R$};
\node[below left] at (0,0) {$-0.5\Delta$};
\node[below] at (7,0) {$+0.5\Delta$};
\node[left] at (0,0.3) {0};
\node[left] at (0,2.0) {1};
\draw[<->] (1,0.3)--(1,2.0);
\node[right] at (1,1.15) {$b_T=0$};
\draw[<->] (5,0.3)--(5,2.0);
\node[right] at (5,1.15) {$b_T=1$};
\node[above] at (3.5,2.0) {翻转点};
\end{tikzpicture}
\caption{dither 扫描下终端 bit 的翻转。32 级均匀 dither 中，$b_T=1$ 的比例
精确反映残差小数部分 $\{R\}$。}
\label{fig:dither}
\end{figure}

数学上，对 dither 电平 $d$，终端决策为 $b_T(d)=\mathbf{1}[R+d>0]$。32 级均匀
dither 在 $[-0.5,+0.5)\Delta$ 上平均后：
\[
\mathbb{E}[b_T] = \frac{R + 0.5\,\Delta}{\Delta} \pmod 1
\]
即 $b_T$ 为 1 的概率等于残差的小数部分。\textbf{dither 把不可见的小数变成
可观测的概率，平均后精确还原。}

\section{dither 电压公式}

\texttt{COM\_ideal.va} L86--L91 实现：
\[
\boxed{v_{dither}(i) = V_{LSB}^{cal}\cdot\!\left(\frac{i+0.5}{N_{lev}}-0.5\right),\quad
i=0,1,\dots,31}
\]
其中 $V_{LSB}^{cal}=0.8274$\,mV $=2(V_{REFP}-V_{REFN})/4351$，$N_{lev}=32$。
序列对称分布在 $[-0.5,+0.5)\cdot V_{LSB}^{cal}$，\textbf{均值恰好为零}，
不引入系统偏差。$(i+0.5)$ 用中点避免取到边界。

\section{时序控制}

三个事件控制 dither 状态：

\begin{enumerate}
\item \textbf{CAL 上升沿}（L67--L72）：\texttt{dither\_index}=0，重启扫描
\item \textbf{CAL 下降沿}（L74--L78）：\texttt{dither}=0，确保正常转换无污染
\item \textbf{CLK 上升沿}（L80--L100）：仅当 \texttt{edge\_in\_pair}==0 时
      更新到下一个 dither 电平，此后 16 个沿（一对 D+/D$-$）保持恒定
\end{enumerate}

\textbf{关键约束}：dither 在一对内保持恒定，使 D+ 和 D$-$ 共用同一 dither。
差分 $R=D_+-D_-$ 时 dither 的\textbf{一阶效应抵消}，残留仅来自 SAR 离散
决策的非线性，由 32 级平均进一步抑制。

%===================================================================
\chapter{SAR 逻辑与定时}
%===================================================================

\section{compare-then-commit 时序}

\texttt{SAR\_LOGIC\_0716.va} 采用"先比较后提交"策略（L13--L29），与传统
"试探-保留/撤销"不同：

\begin{enumerate}
\item PRST 下降沿：所有支路保持 VCM，启动 stage 0，\textbf{不施加 MSB 试探}
\item 第一个比较器周期：观察采样残差的符号
\item CCLK 上升沿：锁存 COMN，\textbf{然后}提交当前支路到对应参考态（L102--L128）
\item 重复直至 stage 13（终端决策，不切电容）
\end{enumerate}

极性约定（\texttt{NORMAL\_REF\_SWAP}=1）：
\begin{itemize}
\item $P>N \Rightarrow$ COMP=1 $\Rightarrow$ BITP=1（负向校正）
\item $P<N \Rightarrow$ COMN=1 $\Rightarrow$ BITN=1（正向校正）
\end{itemize}

\section{SYNC 四相时钟}

\texttt{SYNC\_asnyc.va} 在 CLK00 上升沿后产生 14 个周期，每周期 4 相
（图~\ref{fig:clock}）。默认 4\,ns/决策，14 决策共 56\,ns。

\begin{figure}[ht]\centering
\begin{tikzpicture}[font=\scriptsize]
\draw[->] (0,0)--(9.5,0) node[right]{time (ns)};
\foreach \x/\l in {0/0,2/1,4/2,6/3,8/4}{
  \draw (\x,0.1)--(\x,-0.1); \node[below] at (\x,-0.1){\l};
}
% CLK
\draw[thick,blue] (0,1.0)--(1,1.0)--(1,0.3)--(3,0.3)--(3,1.0)--(4,1.0)
  (4,1.0)--(5,1.0)--(5,0.3)--(7,0.3)--(7,1.0)--(8,1.0);
\node[left] at (0,0.65) {CLK};
\node[above] at (1.5,1.0) {$\uparrow$捕获};
% CCLK
\draw[thick,red] (0,2.0)--(2,2.0)--(2,1.3)--(4,1.3)--(4,2.0)--(6,2.0)
  (6,2.0)--(7,2.0)--(7,1.3)--(9,1.3);
\node[left] at (0,1.65) {CCLK};
\node[above] at (3,2.0) {$\uparrow$提交};
\draw[decorate,decoration={brace,amplitude=3pt}] (0,2.4)--(4,2.4)
  node[midway,above]{1 个决策周期 = 4 ns};
\end{tikzpicture}
\caption{SYNC 四相时钟。CLK 上升沿比较器捕获，CCLK 上升沿 SAR 逻辑提交。}
\label{fig:clock}
\end{figure}

%===================================================================
\chapter{解码器与验证}
%===================================================================

\section{正常解码器}

\texttt{DEC\_CAL\_PHY.va} L596--L698，\textbf{仅在 SOUT 上升沿锁存}一次完整
RAW 码，避免 14 拍中间码泄漏到输出。解码流程：

\begin{enumerate}
\item 收集 14 拍 BP 决策，用校准权重 $w_0$--$w_{12}$ 加权求和得 \texttt{raw\_sum}
\item 冗余偏移消除：$\texttt{centered} = \texttt{raw} - (W_{total}-4095)/2$
      （L662--L680）
\item 若 \texttt{NORMALIZE\_REDUNDANT\_RANGE}=1：比例缩放
      $\texttt{code} = 4095 \times \texttt{raw} / 4351$
\item Q$\to$整数四舍五入，限幅 $[0, 4095]$
\end{enumerate}

\texttt{APPLY\_CAL\_WEIGHT}=1 时用校准后的 $w_0$/$w_1$/$w_2$ 解码，
=0 时用理想 2048/1024/512，用于 A/B 对比。

\section{A/B/C 验证矩阵}

\begin{table}[ht]\centering\footnotesize
\begin{tabular}{p{1.8cm}p{1.2cm}p{1.2cm}p{1.5cm}p{5.0cm}}
\toprule
\textbf{模式} & \textbf{BYP} & \textbf{APL} & \textbf{ENOB} & \textbf{用途} \\
\hline
C（旁路） & 1 & 0 & 12.0 & 纯旁路，理想权重参考 \\
A（不应用） & 0 & 0 & $\approx$12.0 & 校准跑完但解码仍用理想权重 \\
B（应用） & 0 & 1 & 11.5--11.7 & 校准权重进解码器，目标模式 \\
\bottomrule
\end{tabular}
\end{table}

\noindent\footnotesize 注：BYP = \texttt{CAL\_BYPASS}，APL = \texttt{APPLY\_CAL\_WEIGHT}。

\section{诊断规则}

三组对比可定位问题归属：

\begin{itemize}
\item \textbf{C 只有 11.5}：问题在架构或解码器，与校准无关
      $\to$ 检查冗余裕量、端点裁剪
\item \textbf{C=12.0 但 B=11.5}：校准权重引入了误差
      $\to$ 检查 \texttt{weights\_q} 日志是否偏离理想值
\item \textbf{C=12.0, A=12.0, B=11.5}：校准测量误差通过权重应用污染解码
      $\to$ 增加平均次数或检查 calDAC 自失配
\end{itemize}

\section{合并版 ENFORCE\_MIN\_AVG}

\texttt{DEC\_CAL\_PHY.va} 合并版引入 \texttt{ENFORCE\_MIN\_AVG} 参数：
默认 =1（安全），当传入 \texttt{AVG\_PAIRS\_LOG2}=0 时自动 clamp 到 5（32 对）；
设 =0 则允许 1 对快速调试，仿真时间降为 1/32。

%===================================================================
\chapter{性能分析与改进路径}
%===================================================================

\section{ENOB 损耗的定量分析}

实测 ENOB 11.5--11.7 对应 SNDR 71--72\,dB，距 12 位理想 74\,dB 差 2--3\,dB。
四大根因及其定量影响：

\begin{table}[ht]\centering\footnotesize
\begin{tabular}{p{3.2cm}p{1.8cm}p{6.8cm}}
\toprule
\textbf{因素} & \textbf{影响} & \textbf{定量推导} \\
\hline
递归误差传播 & 0.2--0.3 bit &
$w_2$ 误差 $\epsilon_2$ 传播到 $w_1$ 的墙，再传播到 $w_0$。
若 $\epsilon_2=0.1\%$（0.5 LSB），$w_0$ 误差可达 1--2 LSB，在码 2048 边界产生 INL 跳变 \\
冗余裕量过紧 & 偶发 sparkle &
纠错范围 $=256$，需纠正的最大误差 $=255$（stage 3 之后所有位之和），裕量仅 1 LSB。
stage 3/4 同时判错即产生 256 LSB 跳码 \\
端点裁剪 & SFDR 2--3 dB &
\texttt{REMOVE\_REDUNDANCY\_OFFSET} 模式下，总权重 4351 减偏移 128 后范围
$[-128, 4223]$，两端各 128 码饱和，产生 3 次谐波 \\
dither-SAR 耦合 & $\sim$0.1 bit &
同一 dither 影响 7 步 SAR 决策（非独立噪声），若 dither 导致第 1 步判错，
后续 6 步在错误起点收敛，产生系统性偏差，32 对平均后残留 $\sim$0.05\% \\
\bottomrule
\end{tabular}
\end{table}

\section{改进路径}

按实施难度从低到高排列：

\begin{enumerate}
\item \textbf{归一化解码}（改参数）：设 \texttt{NORMALIZE\_REDUNDANT\_RANGE}=1，
      比例缩放消除端点裁剪。预期 SFDR 改善 2--3\,dB

\item \textbf{扩展校准到中位}（换文件）：用
      \texttt{DEC\_CAL\_PHY\_HUANG\_RECURSIVE\_V4.va} 校准 $w_3$/$w_4$/$w_5$
      （权重 256/256R/128），覆盖更多非线性来源。注意 $w_3$ 的墙需用 $w_4$
      单独构建（$R=0$），否则超出 calDAC 范围

\item \textbf{迭代自校准}（改软件）：用已校准的 MSB 反向校准 calDAC 自身权重，
      迭代 2--3 次收敛。消除"用有误差的尺子量东西"问题，不需 DEM 硬件

\item \textbf{DEM}（改硬件）：MSB 分解为单位电容阵列，$\sqrt{N}$ 效应使残余
      误差 $<0.05$ LSB，SNDR 损失 $<0.01$\,dB。代价是 CDAC 与开关阵列重设计
\end{enumerate}

\section{与 Huang 2024 论文对比}

\begin{table}[ht]\centering\footnotesize
\begin{tabular}{p{2.5cm}p{4.5cm}p{5.0cm}}
\toprule
\textbf{维度} & \textbf{Huang 2024} & \textbf{本包 V6} \\
\hline
ADC 规格 & 16-bit, 5\,MS/s, 0.18\,$\mu$m & 12-bit, 行为级 \\
CDAC 权重 & 二进制 + 冗余 & 非二进制 4C+2C 段 \\
校准目标 & 14 个高位 (b6--b19) & \textbf{仅 3 个 MSB} \\
calDAC & 6-bit & 7-bit (stages 6--12) \\
噪声来源 & 真实 $\sigma_{n,tot}$ + SRM & 确定性 dither 32 电平 \\
SRM & 有（22 次额外决策） & 无 \\
Split Sampling & 有（降噪） & 无 \\
2-bit Flash & 有（b18/b19 过范围） & 无 \\
Vos 覆盖 & 对称冗余 $\pm W_5$, 4$\sigma$ & $\pm$35\% 容差门限 \\
实测精度 & SNDR 93.7\,dB, INL $\pm$0.9 LSB & ENOB 11.5--11.7 \\
\bottomrule
\end{tabular}
\end{table}

核心区别：Huang 的 16-bit 设计有完整噪声预算（SRM + Split Sampling 将噪声
从 111\,$\mu$V 降至 38\,$\mu$V），本包作为行为级验证去掉了这些电路技术，
用确定性 dither 替代真实噪声。算法内核（offset-binary + D+/D$-$ + 32 对平均
+ 递归）\textbf{完全一致}。

\section{DEM vs 校准：中小失配的结论}

对于 MSB 2--3 位的中小失配，DEM 在理论上优于前台校准：

\begin{itemize}
\item \textbf{$\sqrt{N}$ 效应}：$w_0=2048$ 分解为 2048 个单位电容，DEM 后残余
      $=0.1\%/\sqrt{2048}=0.0022\%$，即 0.045 LSB，SNDR 损失 $<0.01$\,dB
\item \textbf{不受 calDAC 自失配限制}：DEM 不需要测量工具
\item \textbf{无递归误差传播}：各 MSB 独立轮换
\item \textbf{但要求 CDAC 硬件重设计}为单位电容阵列，当前二进制 CDAC 无法应用
\end{itemize}

推荐路径：CDAC 已定版 $\to$ "MSB 校准 + 中位 DEM"或迭代自校准；
新项目 $\to$ 优先 DEM。

\section{仿真配置}

\begin{itemize}
\item 校准时长：32 对 $\times$ 3 目标 $\approx$ 19.2\,$\mu$s
\item 建议 \texttt{tran stop} $\geq$ 35\,$\mu$s
\item FFT 必须从 \texttt{DONE}=1 后取样
\item 预期日志：\texttt{OFFSAR start} $\to$ 3 条 \texttt{measure}
      $\to$ \texttt{finish done=1 err=0}
      $\to$ \texttt{weights\_q} 接近 $\{131072, 65536, 32768\}$
\end{itemize}

%===================================================================
\appendix
\chapter{核心公式速查}
%===================================================================

\footnotesize
\begin{longtable}{p{4.0cm}p{9.0cm}}
\toprule
\textbf{名称} & \textbf{公式} \\
\hline
\endhead
dither 电压 & $v_d(i)=V_{LSB}^{cal}\!\left(\frac{i+0.5}{N}-0.5\right)$，
              $V_{LSB}^{cal}=\frac{2V_{REF}}{4351}=0.8274$\,mV \\
D+/D$-$ 差分 & $R=D_+-D_-\approx\frac{2(W_t-W_w)}{\Delta}$ \\
Q 格式平均 & $R_q=\mathrm{round}\!\left(64\cdot\frac{\sum(D_+-D_-)}{32}\right)$ \\
测得权重 & $W_{meas}=W_{wall,q}+R_q$ \\
冗余偏移消除 & $\text{centered}=\text{raw}-\frac{W_{total}-4095}{2}$ \\
归一化解码 & $\text{code}=\mathrm{round}\!\left(\frac{4095\cdot\text{raw}}{4351}\right)$ \\
DEM 残余 & $\sigma_{eff}=\frac{\sigma_{unit}}{\sqrt{N}}$ \\
理论分辨率 & $\frac{2}{64\cdot\sqrt{32}}\approx 0.0044$\,LSB $\approx 1.1\times10^{-6}$ \\
\bottomrule
\end{longtable}

\begin{thebibliography}{9}
\bibitem{huang2024} Z. Huang, \textit{Advanced clock multiplier and SAR ADC
  design techniques for high-resolution signal chain systems}, Ph.D. thesis,
  2024.
\end{thebibliography}

\end{document}
